% Generated draft: requires mathematical and source review.
\documentclass[11pt]{article}
\usepackage[margin=25mm]{geometry}
\usepackage{fontspec}
\setmainfont{Latin Modern Roman}
\usepackage{amsmath,amssymb,mathtools,unicode-math}
\setmathfont{Latin Modern Math}
\usepackage{enumitem,needspace,xurl,hyperref}
\setlength{\parindent}{0pt}
\setlength{\parskip}{5pt}
\setlength{\emergencystretch}{3em}
\newcommand{\N}{\mathbb{N}}
\newcommand{\Z}{\mathbb{Z}}
\newcommand{\Q}{\mathbb{Q}}
\newcommand{\R}{\mathbb{R}}
\newcommand{\C}{\mathbb{C}}
\newcommand{\F}{\mathbb{F}}
\newcommand{\A}{\mathbb{A}}
\newcommand{\PP}{\mathbb{P}}
\newcommand{\E}{\mathbb{E}}
\newcommand{\eps}{\varepsilon}
\newcommand{\dd}{\,\mathrm{d}}
\begin{document}
\section*{UnsolvedMath ID: 20002818}
\textbf{Source problem code:} catalogue entry
\subsection{Definitions and mathematical statement}
Let $J$ be a finite-dimensional nilpotent associative algebra over $\F_q$, $G=1+J$, and fix a nontrivial additive character $\vartheta:\F_q\to\C^\times$. For $\lambda\in J^*$ write $\mathcal O_\lambda=G\lambda G$ for its two-sided dual orbit, $d_\lambda=|G\lambda|$ for its left-orbit size, and use the standard algebra-group normalization
\[
 \chi_\lambda(1+x)
   =c_\lambda\sum_{\mu\in\mathcal O_\lambda}\vartheta(\mu(x)),
 \qquad c_\lambda=d_\lambda/|\mathcal O_\lambda|.
\]
The source asks when every supercharacter is a pointwise product of supercharacters that are themselves irreducible ordinary characters. It recalls the property for $U_n(q)$, says it fails for some pattern groups, and asks for a criterion even for algebra groups. The normalization is part of this question.

The supplied literature assessment attributes to Marberg a positive family: for every pattern subgroup normal in $U_n(q)$, all standard supercharacters factor into irreducible elementary supercharacters. The exact reference is not supplied. The aid then gives a complete commutative-algebra criterion. If $J$ is commutative, its one- and two-sided dual orbits coincide. A standard supercharacter is irreducible exactly when
\[
 \lambda(J^2)=0.
\]
In that case its orbit is a singleton and $\chi_\lambda(1+x)=\vartheta(\lambda(x))$ is linear. Products of these linear standard supercharacters remain of that same form, so
\[
 \chi_\lambda\text{ factors into irreducible standard supercharacters}
 \quad\Longleftrightarrow\quad\lambda(J^2)=0.
\]
Consequently
\[
 \text{every standard supercharacter of }1+J\text{ factors}
 \quad\Longleftrightarrow\quad J^2=0
 \qquad(J\text{ commutative}).
\]
Although the group $1+J$ is abelian in this case, its standard supercharacter theory need not be its minimal character theory; reducible standard supercharacters are not rescued by the existence of other linear ordinary characters.

For general $J$, the aid's finite orbit test is normalization sensitive. Given prospective irreducible standard factors $\chi_{\lambda_1},\ldots,\chi_{\lambda_r}$, define
\[
 N(\nu)=\#\{(\mu_1,\ldots,\mu_r)\in
           \mathcal O_{\lambda_1}\times\cdots\times\mathcal O_{\lambda_r}:
                 \mu_1+\cdots+\mu_r=\nu\}.
\]
Additive Fourier independence gives the exact equivalence
\[
 \chi_\lambda=\prod_{j=1}^r\chi_{\lambda_j}
 \quad\Longleftrightarrow\quad
 \left(\prod_jc_{\lambda_j}\right)N(\nu)
           =c_\lambda\mathbf1_{\mathcal O_\lambda}(\nu)
 \quad\text{for every }\nu\in J^*.
\]
Equivalently, the orbit-addition map has precisely $\mathcal O_\lambda$ as image, has uniform fibers, and the degrees satisfy $d_\lambda=\prod_jd_{\lambda_j}$. Irreducibility of a standard factor can itself be tested by $|\mathcal O_{\lambda_j}|=d_{\lambda_j}^2$, from its character norm. Uniform support alone without the scalar normalization or degree condition is insufficient. This finite criterion and the commutative classification do not classify all noncommutative algebra or pattern groups with the property.
\subsection{Short English statement}
Classify standard supercharacters that factor into irreducible standard supercharacters. For commutative algebra groups this happens exactly for functionals annihilating $J^2$, so the property holds for all supercharacters exactly when $J^2=0$. In general the aid gives a normalized uniform-fiber orbit-addition test.
\subsection{Sources}
\begin{itemize}
\item American Institute of Mathematics workshop \emph{Supercharacters and combinatorial Hopf algebras}, Problem 1.6, \url{https://aimath.org/WWN/supercharacters/supercharacters.pdf}.
\item Supplied assessment attributes the normal-pattern-subgroup factorization theorem to Marberg, without title, year or complete bibliographic details.
\item Local UnsolvedMath registry, permanent ID 20002818, source code \texttt{catalogue entry}, canonical location \texttt{aim-representation-theory-notes.json}, \texttt{notes[89]}, record dated 14 August 2026. The original literature field is empty.
\item Supplied research attempt 1: partial result, problem status at run partially solved, rigor labelled proved, novelty confidence low. Its commutative classification and orbit-fiber test are machine-generated research aid, not peer reviewed; independently verify before citation.
\end{itemize}
\end{document}
